% 正控制样本 6（判据全绿）：可选线宽参数 —— 三条主线用 \toprule[..]/\midrule[..]/\bottomrule[..]，
% 分组线用 \cmidrule[..](trim){range}。用途 = I-1「第三处假红」的**防假红对照**（2026-10-03）：
% 旧 RULES_RE 不吞 [..] ⇒ 剥掉规则线后留裸 `[..]` ⇒ 幻影一行 ⇒ C5 + C8 假红（表本身能编过）。
% ★ 紧贴（\toprule[1.2pt]）与**横空白分隔**（\midrule [0.8pt]）两种写法都覆盖（不声称穷尽）。
% 本件必须全绿。列数 = 3；行数 = 4（分组头行 / 列头行 / 两条数据行）。
% mcm-table-src: r1c1 <- pasted:r1c1
% mcm-table-src: r1c2 <- pasted:r1c2
% mcm-table-src: r1c3 <- pasted:r1c3
% mcm-table-src: r2c1 <- pasted:r2c1
% mcm-table-src: r2c2 <- pasted:r2c2
% mcm-table-src: r2c3 <- pasted:r2c3
% mcm-table-src: r3c1 <- pasted:r3c1
% mcm-table-src: r3c2 <- pasted:r3c2
% mcm-table-src: r3c3 <- pasted:r3c3
% mcm-table-src: r4c1 <- pasted:r4c1
% mcm-table-src: r4c2 <- pasted:r4c2
% mcm-table-src: r4c3 <- pasted:r4c3
\begin{table}[htbp]
  \centering
  \caption{A compliant table with explicit rule thicknesses.}
  \begin{tabular}{l S[table-format=1.2] S[table-format=1.2]}
    \toprule[1.2pt]
    & \multicolumn{2}{c}{Score} \\
    \cmidrule[0.6pt](lr){2-3}
    Item & {A} & {B} \\
    \midrule [0.8pt]
    Alpha & 1.00 & 2.00 \\
    Beta  & {--} & 3.00 \\
    \bottomrule[1.2pt]
  \end{tabular}
\end{table}
